\documentclass[12pt,aspectratio=169]{beamer}
\usepackage[T1]{fontenc}
\usepackage{array,tabularx}

\definecolor{GEInk}{HTML}{132A3E}
\definecolor{GETeal}{HTML}{17675E}
\definecolor{GEFooter}{HTML}{EAF2F1}

\setbeamercolor{background canvas}{bg=white}
\setbeamercolor{normal text}{fg=GEInk,bg=white}
\setbeamercolor{structure}{fg=GETeal}
\setbeamercolor{frametitle}{fg=GEInk}
\setbeamercolor{GE footer}{fg=GEInk,bg=GEFooter}
\setbeamerfont{title}{series=\bfseries}
\setbeamerfont{frametitle}{series=\bfseries}
\setbeamerfont{GE footer}{size=\scriptsize}
\setbeamersize{text margin left=9mm,text margin right=9mm}
\setbeamertemplate{navigation symbols}{}
\setbeamertemplate{frametitle}{%
  \vskip0.5em
  \strut\insertframetitle\strut\par
  \ifx\insertframesubtitle\empty\else
    {\usebeamerfont{framesubtitle}\usebeamercolor[fg]{framesubtitle}%
     \strut\insertframesubtitle\strut\par}%
  \fi
}
\setbeamertemplate{footline}{%
  \begin{beamercolorbox}[wd=\paperwidth,ht=2.7ex,dp=1.1ex,
    leftskip=9mm,rightskip=9mm]{GE footer}%
    \usebeamerfont{GE footer}\insertsectionhead\hfill
    \insertframenumber\,/\,\inserttotalframenumber
  \end{beamercolorbox}%
}

% Use short text or math arguments. Put surrounding operators outside the box.
\makeatletter
\newcommand{\GEHighlightBox}[2]{{%
  \setlength{\fboxsep}{1.5pt}%
  \colorbox{#1}{\color{GEInk}#2}%
}}
\newcommand{\GEHighlightMath}[3]{\GEHighlightBox{#1}{$\m@th#2#3$}}
\newcommand{\GEHighlight}[2]{%
  \ifmmode
    \mathpalette{\GEHighlightMath{#1}}{#2}%
  \else
    \GEHighlightBox{#1}{#2}%
  \fi
}
\makeatother

% Replace these illustrative roles, metadata, and frames for each lesson.
\definecolor{GERevenue}{HTML}{D8EFE9}
\definecolor{GECost}{HTML}{F7E8D2}
\title{Revenue, cost, and profit}
\subtitle{Global Economics}
\date{}

\begin{document}

\section{Learning objective}
\begin{frame}
  \titlepage
  \centering
  \textbf{Learning objective}\par\medskip
  Distinguish revenue from cost and use both to calculate profit.
\end{frame}

\section{A market stall}
\begin{frame}[t]{How much did the stall earn after costs?}
  In this example, a stall takes in 90 thousand COP from sales and
  spends 60 thousand COP in total. Which amount is revenue, and how much
  is left after costs?

  \medskip
  \renewcommand{\arraystretch}{1.25}
  \begin{tabularx}{\linewidth}{@{}>{\raggedright\arraybackslash}X
                                     >{\raggedright\arraybackslash}X@{}}
    \textbf{\GEHighlight{GERevenue}{Revenue}} &
    \textbf{\GEHighlight{GECost}{Total cost}} \\
    \hline
    \visible<2->{Money received from sales} &
    \visible<2->{Money spent in total} \\
    \visible<2->{\GEHighlight{GERevenue}{90 thousand COP}} &
    \visible<2->{\GEHighlight{GECost}{60 thousand COP}}
  \end{tabularx}

  \visible<3->{%
    \[
      \text{Profit (thousand COP)}
        = \GEHighlight{GERevenue}{90} - \GEHighlight{GECost}{60} = 30
    \]
  }
\end{frame}

\section{Summary}
\begin{frame}[t]{What connects the three amounts?}
  \[
    \text{Profit} = \GEHighlight{GERevenue}{\text{Revenue}}
                   - \GEHighlight{GECost}{\text{Total cost}}
  \]
  Express all three amounts in the same units for the same period.
\end{frame}

\end{document}
